% =====================================================================
% A2 — Section 6: Discussion and conclusion.  DRAFT 1, 6 Oct 2026.
% The one "earned" contrast sentence (determine vs. find) lives here and in
% the introduction only. The mechanism is stated as an interpretation.
% =====================================================================
\section{Discussion and conclusion}
\label{sec:discussion}

The first part of the paper's argument holds up well. Whether a training set
pins down the linear representation can be computed exactly, and both my implementation and the authors' give a transition at $r \approx 0.4$. The second part holds
up less well. With the released settings, trained networks find the structure less reliably than the theory suggests, and the time they need does not
follow $1/\lambda_3$.

In terms of Fig.~\ref{fig:pipeline}, the effective theory also predicts the
lower half, but only for an gradient flow without a decoder. It says when the data \emph{determine} the representation, and not if a trained
network \emph{finds} it. One explanation that fits my results is that the
decoder sometimes overfits the training set before the embeddings have become
structured. Once that happened, the training loss is close to zero. 
This would explain why most failed runs stop improving. The effective loss has no decoder, so it cannot
describe this. Slowing the decoder down helps and speeding it up makes it worse, which
matches the competition between representation learning and decoder
overfitting that the paper describes in its Sec.~4. But I did not measure gradients directly.

My results show that C1 is correct, and
the published Fig.~4(b) may come from settings closer to the slow-decoder runs.
They also show that the agreement in that figure depends on
training settings the paper does not report, and that $\lambda_3 > 0$ alone does not predict whether a network will find the
structure.

The study has some limitations. I only looked at addition with $p = 10$ and
one-dimensional embeddings, and since training uses AdamW rather than gradient
flow, only trends can be compared with the theory. With twenty seeds per size,
the success probabilities at intermediate fractions are still uncertain by
about $\pm 0.2$.

For reproducibility, it would have helped if the paper had stated the
settings and the number of seeds behind each figure. Reporting runs that never
succeed as a separate group would also have helped. Both turned out to matter
for me to understand what the published figures show.
